\begin{studiumworked}{M06. Dos raíces y primer bloqueo · Intermedio}
\textbf{Enunciado.} Una tobera ideal tiene $A_e/A_t=2.86$ y $\gamma=1.4$. Calcule las dos raíces de salida para garganta sónica y la contrapresión relativa de primer bloqueo. Compare con una convergente simple.

\textbf{Datos.} Expansión: $2.86$; Exponente: $1.4$.

\textbf{Incógnitas.} Magnitudes solicitadas y cierre físico.

\textbf{Hipótesis.} Modelo académico con hipótesis del capítulo. Datos originales; no representan hardware ni ensayo.

\textbf{Modelo físico.} Conservación y relaciones constitutivas declaradas.
\[A/A^*=f(M,\gamma)\]
\[(p_b/p_0)_{inicio}=[1+(\gamma-1)M_{e,sub}^2/2]^{-\gamma/(\gamma-1)}\]
\textbf{Desarrollo.}
\begin{enumerate}
\item La inversión se encierra en dos intervalos distintos: $(0,1)$ y $(1,\infty)$. Las raíces son $0.2076227$ y $2.586722$, respectivamente.
\item El primer bloqueo usa la raíz subsónica del divergente. Sustituirla da $p_b/p_0=0.9704012$. Para una convergente simple el umbral sería $0.5282818$.
\item La diferencia aparece porque el divergente subsónico recupera presión después del mínimo sónico. Ninguno de esos valores es automáticamente la presión adaptada de la rama supersónica.
\end{enumerate}
\textbf{Resultados.} $M_{sub}=0.2076227$; $M_{sup}=2.586722$; $(p_b/p_0)_{inicio}=0.9704012$.

\textbf{Interpretación.} El resultado corresponde a las hipótesis indicadas; su interpretación requiere conservar fronteras y unidades.

\textbf{Comprobación.} Sustituir cada raíz en la relación área-Mach recupera la misma expansión dentro de la tolerancia numérica.

\textbf{Límites.} No se calcula separación ni posición de choque; el umbral es inviscido y cuasiunidimensional.
\end{studiumworked}
